\section{Conclusiones del segundo corte}

\begin{enumerate}

\item El paso del diseño conceptual al prototipo obligó a ajustar varias decisiones del primer corte a los componentes realmente disponibles: el circuito de paso único hacia el tanque de efluentes, la bomba periférica como bomba de impulsión, la agitación con una bomba sumergible, dos SSR sin contactores, un flotador doble para el nivel y la supervisión remota con una Raspberry Pi. Estos cambios se registran en la Sección~\ref{sec:cambios_prototipo}, y el P\&ID, los requisitos y el plan de validación se actualizaron para que describan el prototipo construido y no la propuesta inicial.

\item Se construyó el prototipo físico sobre dos carros (Sección~\ref{sec:montaje}). En el tanque de mezcla de plástico se instalaron la resistencia de inmersión, el sensor de temperatura, el flotador doble y la bomba de agitación; además se instalaron la bomba de impulsión, la electroválvula de llenado XV-102 y un tablero que separa la zona de potencia del banco del PLC (Sección~\ref{sec:montaje_instrumentacion}). Quedan pendientes de instalación la bomba dosificadora P-101 y la electroválvula de descarga XV-103, que aún no se ha comprado y será de solenoide de acción directa, porque la descarga desde el tanque de mezcla ocurre por gravedad. Los datos de placa de las bombas fijan límites que deben respetarse en la operación: 40\,$^{\circ}$C para la bomba de impulsión y 35\,$^{\circ}$C para la bomba de agitación.

\item La secuencia del ciclo CIP se diseñó con un GRAFCET y se programó en el PLC en una versión preliminar (Ladder v1.2) con las tres fases del ciclo: preenjuague, lavado con detergente y enjuague final (Sección~\ref{sec:montaje_programacion}). Sus tiempos son preliminares, y aún faltan la pausa para el vaciado del tanque de efluentes y el estado de parada segura.

\item El ensayo en lazo abierto con escalones de potencia permitió identificar experimentalmente el modelo térmico del tanque de mezcla, con una constante de tiempo de 88,4\,min y una ganancia de 0,0733\,$^{\circ}$C/W (Ecuación~\eqref{eq:modelo_experimental}). Con este modelo se diseñó un controlador PI por lugar de las raíces, afinado por simulación con saturación, retardo y anti-windup ($K_p = 12\,\%/^{\circ}\mathrm{C}$, $T_i = 1800\,\mathrm{s}$), que en simulación tiene un sobreimpulso de 0,58\,°C en el peor caso de la envolvente con anti-windup. Se programó en el PLC con el bloque \texttt{PID\_Compact} y modulación de la potencia por proporción de tiempo, pero todavía no se ha probado en lazo cerrado (Sección~\ref{sec:diseno_pi}).

\item La cadena de supervisión remota quedó implementada: la Raspberry Pi lee las variables del PLC con Node-RED mediante el protocolo S7 y las envía a una API en AWS, que se consulta desde una interfaz en Streamlit. El control permanece en el PLC y no depende de esta capa (Sección~\ref{sec:montaje_comunicaciones}).

\item La validación del prototipo frente a los requisitos sigue incompleta. El ensayo térmico y el registro de datos del PLC dieron la base experimental del control, pero quedan pendientes las pruebas de la secuencia, de los niveles, de la dosificación, de las protecciones y del controlador en lazo cerrado (Sección~\ref{sec:montaje_pruebas}).

\end{enumerate}

\section{Trabajo previsto para el tercer corte}
\noindent Durante el tercer corte el proyecto pasará del prototipo construido a su validación experimental frente a los requisitos.
\par\medskip
\noindent\textbf{Completar el prototipo:} Se instalará la bomba dosificadora P-101 con su convertidor de 12\,V\,DC, se comprarán e instalarán una electroválvula de descarga XV-103 de acción directa y la válvula manual de vaciado HV-101, y se terminarán y documentarán el conexionado del paro de emergencia y el diagrama de conexiones del PLC.
\par\medskip
\noindent\textbf{Programación de la secuencia:} Se implementarán en el PLC la pausa para el vaciado manual del tanque de efluentes entre fases y el estado de parada segura. Se ajustarán los tiempos de las fases y se corregirá el GRAFCET según la norma IEC 60848.
\par\medskip
\noindent\textbf{Control de temperatura:} Se probará el controlador PI en lazo cerrado sobre el prototipo, se comparará su respuesta con la simulada y, si es necesario, se reajustarán sus parámetros. Después se integrará en la secuencia del ciclo, con una consigna compatible con los límites de temperatura de las bombas. También se corregirán los dos ajustes identificados en la máquina de estados del registro de datos.
\par\medskip
\noindent\textbf{Supervisión:} Se confirmarán las direcciones IP definitivas y el periodo de actualización de los datos, y se documentará la evidencia de funcionamiento de la cadena de supervisión remota.
\par\medskip
\noindent\textbf{Validación experimental:} Se ejecutarán las pruebas del plan de validación (Sección~\ref{sec:plan_validacion}) y se registrará para cada requisito si se cumple. Se definirá además el método para evaluar la efectividad de la limpieza de la tubería de prueba.
\par\medskip
\noindent\textbf{Documentación:} Se cerrará el informe técnico con los resultados, se completará el artículo en inglés y se grabará el video del proceso.
